% ----------------------
\subsection{Section 91 (9 March 1833)}
% ----------------------
	\label{DC91}
	
	% -----
	\subsubsection{Historical Background}
	% -----
		\label{DC91-HB}
		
		% --
		\paragraph{JSP}
		% --
		
			``At the time this revelation was dictated, JS was actively engaged in revising the Bible. A month earlier, on 2 February, scribe Frederick G. Williams noted, `This day completed the translation and the reviewing of the New testament and sealed [it] up no more to be broken till it goes to Zion.' A revelation dated 8 March, one day earlier than the date of the revelation featured here, referenced JS's continued work on several books in the latter part of the Old Testament and made JS's other duties contingent upon his completion of the Bible revision: `Now verely I say unto you I give unto you a commandment that you continue in this ministry and presidency and when you have finished the translation of the prophets you shall from thence forth preside over the affairs of the Church and the School.' Emphasizing the importance of completing his scriptural revisions, this declaration may have led JS to also ponder at what point the work would be finished and whether the `translation of the prophets' included work on the books of the Apocrypha.
			
			``The word \textit{apocrypha} transliterates a Greek term meaning `hidden' or `concealed.' Protestants used the term to denote the dozen or so books not found in the Hebrew canon but that were included in the Greek version of the Old Testament known as the Septuagint. Because the Septuagint was the basis for early Latin translations of the Old Testament, the Apocrypha was integrated into early Christian Bibles. Although the Catholic scholar Jerome hesitated to include the books along with his translation of the Hebrew canon in the fourth century, they eventually came to be widely regarded as scriptural in medieval Christianity. Martin Luther, in his 1534 translation of the Bible, however, placed the books of the Apocrypha at the end of the Old Testament and explained that they were `not held equal to the Holy Scriptures, and yet are profitable and good to read.' Throughout the Reformation, subsequent editions of the Bible also placed these books in a separate section labeled `Apocrypha.' The Roman Catholic Church responded by affirming the canonical status of all but three of the books at the Council of Trent in 1546. Thus, acceptance or rejection of the authority of the Apocrypha was at times part of the larger religious dispute surrounding the Reformation.
			
			``Shortly before JS began his translation of the Book of Mormon, a widely publicized and acrimonious debate took place within the British and Foreign Bible Society over the inclusion of the Apocryphal books within the Bibles it distributed. Some members of the society accused their leaders of `adulterating the Scriptures, by circulating the lies and fables of the \textit{Apocrypha} along with the words of eternal life.' After multiple schisms, the society finally relented in 1826--1827 and adopted a resolution banning the distribution of Bibles containing the Apocrypha. By April 1828, the American Bible Society, witnessing the fallout caused by arguments within its British counterpart, made the same decision. Public disputes about the Apocrypha nevertheless continued in the United States. In 1832, an Ohio newspaper railed against the Apocrypha in an article designed to prove the veracity of the accepted canonical books of the Old and New Testaments; the article argued that `the claims of the Apocrypha must be despatched in few words' and that the authors of it could not `comport with a claim to a divine origin.' In his widely published 1833 letters against Catholicism, Calvinist theologian William Craig Brownlee specifically denounced his rivals who `decorate the apocrypha with the honors of inspiration.'
			
			``Despite the general movement of American Protestants away from the Apocrypha, early Mormons did not seem to view the Apocrypha with the same disdain, perhaps owing to their foundational belief in the existence of scriptural books outside of the Old and New Testaments, such as the Book of Mormon. An April 1829 revelation explained that there were `records which contain much of my gospel, which have been kept back because of the wickedness of the people' and that `parts of my scriptures . . . have been hidden because of iniquity.' The Book of Mormon also points to the existence of sacred texts outside the traditional biblical canon, asserting that `many parts which are plain and most precious' were lost from the writings that eventually constituted the Bible. Early members of the church apparently discussed regularly the notion that teachings were missing from the traditional canon. JS's history, for instance, indicates that concern for missing biblical books served as the background for his lengthy dictation in December 1830 termed the `Prophecy of Enoch,' which provided additional information about early biblical prophets:
			
			\begin{quote}
				Much conjecture and conversation frequently occurred among the saints, concerning the books mentioned and referred to, in various places in the old and new testaments, which were now no where to be found. The common remark was, they are `lost books'; but it seems the apostolic churches had some of these writings, as Jude mentions or quotes the prophecy of Enoch the seventh from Adam. To the joy of the little flock, which in all, from Colesville to Canandaigua, numbered about seventy members, did the Lord reveal the following doings of olden time from the prophecy of Enoch.
			\end{quote}
			
			``The concern over lost books of scripture apparently continued after 1830; JS, for instance, wrote to church members in Missouri in June 1833, saying, `We have not found the book of Jasher nor any of the othe[r] lost books mentioned in the bible as yet nor wille we obtain them at present.' This acceptance of the notion of additional scripture apparently extended to the Apocrypha in the minds of some early Mormons. In July 1832, the first church-owned newspaper quoted from 2 Maccabees and seemed to deride the treatment of the Apocrypha as noncanonical, explaining that it was `the wisdom of man' that `has seen fit to call [it] Apocrypha.' In January 1833, the same newspaper again derided negative attitudes toward Apocryphal works when it published an excerpt from the Apocryphal addition to the book of Esther, `which the ancient men of the world, put down as doubtful.'
			
			``The revelation featured here affirmed `there are many things contained' in the Apocrypha `that are true.' Nevertheless, this revelation instructed JS that he need not translate the Apocrypha along with the other, canonical books of the Bible, and he apparently never did. He repeated the teachings of the revelation in a letter to the Missouri members a few months later, telling them that `respecting the Apochraphy the Lord Said to us that there were many things in it which were true and ther were many things in it which were not true and to those who desired, it should be given by the spirit to know true from the false.'''
			
	% -----
	\subsubsection{Versions}
	% -----
		\label{DC91-V}
		
		See the \href{https://drive.google.com/file/d/1goAvNz8jS4R8slYfMG_db55Ty2z_vmgK/view?usp=sharing}{sections comparison} document.
	
		\begin{compactdesc}
			\item[JSP] \hyperref[EMS-1833-JS-9-MAR-1833]{9 March 1833}
			\item[BC] ---
			\item[DC 1835] Section 92, 231
			\item[DC 1844] Section 93, 359
			\item[TS] 5:23 (15 December 1844), 737	
			\item[MS] 14:24 (7 August 1852), 376
		\end{compactdesc}

	% -----
	\subsubsection{Section 91:1} % *DC91-1 
	% -----
		\label{DC91-1}

		VERILY, thus saith the Lord unto you concerning the Apocrypha---There are many things contained therein that are true, and it is mostly translated correctly;

		% \begin{framed}
		% \scriptsize
			% \begin{compactdesc}
				% \item[JSP] 
				% % \item[BC] 
				% % \item[]
			% \end{compactdesc}
		% \end{framed}

	% -----
	\subsubsection{Section 91:2} % *DC91-2 
	% -----
		\label{DC91-2}

		There are many things contained therein that are not true, which are interpolations by the hands of men.

		% \begin{framed}
		% \scriptsize
			% \begin{compactdesc}
				% \item[JSP] 
				% % \item[BC] 
				% % \item[]
			% \end{compactdesc}
		% \end{framed}

	% -----
	\subsubsection{Section 91:3} % *DC91-3 
	% -----
		\label{DC91-3}

		Verily, I say unto you, that it is not needful that the Apocrypha should be translated.

		% \begin{framed}
		% \scriptsize
			% \begin{compactdesc}
				% \item[JSP] 
				% % \item[BC] 
				% % \item[]
			% \end{compactdesc}
		% \end{framed}

	% -----
	\subsubsection{Section 91:4} % *DC91-4 
	% -----
		\label{DC91-4}

		Therefore, whoso readeth it, let him understand, for the Spirit manifesteth truth;

		% \begin{framed}
		% \scriptsize
			% \begin{compactdesc}
				% \item[JSP] 
				% % \item[BC] 
				% % \item[]
			% \end{compactdesc}
		% \end{framed}

	% -----
	\subsubsection{Section 91:5} % *DC91-5 
	% -----
		\label{DC91-5}

		And whoso is enlightened by the Spirit shall obtain benefit therefrom;

		% \begin{framed}
		% \scriptsize
			% \begin{compactdesc}
				% \item[JSP] 
				% % \item[BC] 
				% % \item[]
			% \end{compactdesc}
		% \end{framed}

	% -----
	\subsubsection{Section 91:6} % *DC91-6 
	% -----
		\label{DC91-6}

		And whoso receiveth not by the Spirit, cannot be benefited. Therefore%
			\footnote{%
				% \label{}%
				I don't totally see how this ``therefore'' comes out of the previous verses, or of the previous sentence of this verse.
				}
		it is not needful that it should be translated. Amen.

		% \begin{framed}
		% \scriptsize
			% \begin{compactdesc}
				% \item[JSP] 
				% % \item[BC] 
				% % \item[]
			% \end{compactdesc}
		% \end{framed}